\section{结构化分析}

结构化分析（Structured Analysis，SA）是软件工程中一种传统的需求分析与系统设计方法，兴起于20世纪70年代，主要应用于瀑布模型开发流程。它以数据流为核心，通过分层分解的方式描述系统功能。

\subsection{结构化分析的优势}

\subsubsection{需求可视化}

通过图形化建模（如DFD的0层、1层分解）帮助非技术人员理解复杂业务流程。

示例：医院挂号系统中，"患者提交请求→系统验证→生成挂号单"的流程清晰可见。

\subsubsection{模块化设计}

将系统分解为"订单处理""支付验证"等独立模块，降低开发复杂度。适用于团队分工明确的瀑布模型，例如大型电信计费系统开发。

\subsubsection{文档标准化}

生成严格的需求规格说明书（SRS），符合ISO/IEC标准，适合合同约束型项目。

\subsection{结构化分析的局限性}

\subsubsection{需求变更适应性差}

依赖前期完整的需求定义，若在开发后期出现需求变更（如新增"订单退款状态"），需重构多层DFD和数据字典，维护成本高。

对比：敏捷方法通过用户故事和迭代可快速响应变化。

\subsubsection{面向数据流，忽略对象交互}

适合数据处理密集型系统（如银行交易系统），但难以表达复杂的业务规则和对象间的动态交互。

现代替代：UML用例图和时序图更适合面向对象系统的交互建模。

\subsection{自顶向下逐层分解}

结构化分析采用自顶向下的分解策略，从系统的整体功能出发，逐层细化到具体的功能模块。这种方法有助于理解系统的层次结构，确保分析的完整性和一致性。分解过程中需要明确每个层次的功能边界和接口关系，避免功能重叠或遗漏。

\subsection{系统流程图与数据流图}

系统流程图描述了数据在系统中的处理流程，展示了系统的控制逻辑和决策点。数据流图（DFD）则更加关注数据的流动和变换，通过图形符号清晰地表示数据存储、处理过程和外部实体之间的关系。DFD分为不同层次，从顶层的环境图到详细的处理规格说明，形成了完整的系统功能描述体系。

\subsection{数据字典的编制}

数据字典是结构化分析的重要成果之一，它详细定义了系统中所有数据元素的名称、类型、长度、取值范围和业务含义。数据字典为数据流图提供了精确的语义支持，确保不同层次图形之间的一致性。良好的数据字典不仅是分析阶段的重要文档，也是后续设计和实现阶段的重要参考。

\subsection{现代应用与发展}

虽然结构化分析方法在现代软件开发中的主导地位已被面向对象和敏捷方法所取代，但其核心思想仍然具有重要价值。在某些特定领域，如数据处理密集型系统、批处理系统等，结构化分析仍然是有效的分析工具。同时，结构化分析的分解思想和建模方法也被现代方法所吸收和发展，形成了更加灵活和实用的分析技术。